\documentclass{amsart}
% For dealing with references we use the comment environment
\usepackage{verbatim}
\newenvironment{reference}{\comment}{\endcomment}
%\newenvironment{reference}{}{}
\newenvironment{slogan}{\comment}{\endcomment}
\newenvironment{history}{\comment}{\endcomment}

% For commutative diagrams we use Xy-pic
\usepackage[all]{xy}

% We use 2cell for 2-commutative diagrams.
\xyoption{2cell}
\UseAllTwocells

% We use multicol for the list of chapters between chapters
\usepackage{multicol}

% This is generally recommended for better output
\usepackage{lmodern}
\usepackage[T1]{fontenc}

% For cross-file-references

% Package for hypertext links:
\usepackage{hyperref}

% For any local file, say "hello.tex" you want to link to please

% Theorem environments.
%
\theoremstyle{plain}
\newtheorem{theorem}[subsection]{Theorem}
\newtheorem{proposition}[subsection]{Proposition}
\newtheorem{lemma}[subsection]{Lemma}

\theoremstyle{definition}
\newtheorem{definition}[subsection]{Definition}
\newtheorem{example}[subsection]{Example}
\newtheorem{exercise}[subsection]{Exercise}
\newtheorem{situation}[subsection]{Situation}

\theoremstyle{remark}
\newtheorem{remark}[subsection]{Remark}
\newtheorem{remarks}[subsection]{Remarks}

\numberwithin{equation}{subsection}

% Macros
%
\def\lim{\mathop{\mathrm{lim}}\nolimits}
\def\colim{\mathop{\mathrm{colim}}\nolimits}
\def\Spec{\mathop{\mathrm{Spec}}}
\def\Hom{\mathop{\mathrm{Hom}}\nolimits}
\def\Ext{\mathop{\mathrm{Ext}}\nolimits}
\def\SheafHom{\mathop{\mathcal{H}\!\mathit{om}}\nolimits}
\def\SheafExt{\mathop{\mathcal{E}\!\mathit{xt}}\nolimits}
\def\Sch{\mathit{Sch}}
\def\Mor{\mathop{\mathrm{Mor}}\nolimits}
\def\Ob{\mathop{\mathrm{Ob}}\nolimits}
\def\Sh{\mathop{\mathit{Sh}}\nolimits}
\def\NL{\mathop{N\!L}\nolimits}
\def\CH{\mathop{\mathrm{CH}}\nolimits}
\def\proetale{{pro\text{-}\acute{e}tale}}
\def\etale{{\acute{e}tale}}
\def\QCoh{\mathit{QCoh}}
\def\Ker{\mathop{\mathrm{Ker}}}
\def\Im{\mathop{\mathrm{Im}}}
\def\Coker{\mathop{\mathrm{Coker}}}
\def\Coim{\mathop{\mathrm{Coim}}}

% Boxtimes
%
\DeclareMathSymbol{\boxtimes}{\mathbin}{AMSa}{"02}

%
% Macros for moduli stacks/spaces
%
\def\QCohstack{\mathcal{QC}\!\mathit{oh}}
\def\Cohstack{\mathcal{C}\!\mathit{oh}}
\def\Spacesstack{\mathcal{S}\!\mathit{paces}}
\def\Quotfunctor{\mathrm{Quot}}
\def\Hilbfunctor{\mathrm{Hilb}}
\def\Curvesstack{\mathcal{C}\!\mathit{urves}}
\def\Polarizedstack{\mathcal{P}\!\mathit{olarized}}
\def\Complexesstack{\mathcal{C}\!\mathit{omplexes}}
% \Pic is the operator that assigns to X its picard group, usage \Pic(X)
% \Picardstack_{X/B} denotes the Picard stack of X over B
% \Picardfunctor_{X/B} denotes the Picard functor of X over B
\def\Pic{\mathop{\mathrm{Pic}}\nolimits}
\def\Picardstack{\mathcal{P}\!\mathit{ic}}
\def\Picardfunctor{\mathrm{Pic}}
\def\Deformationcategory{\mathcal{D}\!\mathit{ef}}

\setlength{\emergencystretch}{3em}
\begin{document}
\title[Stacks Project, Tag 09DB]{394 --- The first Tor defect of the cotangent complex of a tensor product}
\maketitle
\noindent Copyright (C) 2005--2025 Johan de Jong. Stacks Project authors. Source: \href{https://stacks.math.columbia.edu/tag/09DB}{Tag 09DB}.

\noindent Editorial difficulty note (not author text). Difficulty H2: The proof computes the derived failure of cotangent base change using degree-two Lichtenbaum-Schlessinger relations. It compares relation modules after tensoring, quotients out trivial relations, and identifies the cokernel as Tor rather than merely obtaining a vanishing bound. The required compatibility between the explicit complexes and derived tensoring is central..

\noindent Editorial provenance note (not author text). History: excerpt from revision \texttt{a0444\allowbreak 6e57e\allowbreak c1fbc\allowbreak 252a8\allowbreak 71afc\allowbreak ec775\allowbreak 2fb28\allowbreak 07b14}, cotangent.tex, lines 2846--2911. Reference presentation was adapted; original-author context and core support blocks were retained or added. The mathematical wording is unchanged. Licensed under GNU FDL 1.2 or later; no invariant sections or cover texts. See ../licenses/COPYING. Context blocks reproduce author definitions and supporting results; they are not additional counted problems.

\begin{definition}
\label{definition-cotangent-complex-ring-map}
The {\it cotangent complex} $L_{B/A}$ of a ring map $A \to B$
is the complex of $B$-modules associated to the simplicial $B$-module
$$
\Omega_{P_\bullet/A} \otimes_{P_\bullet, \epsilon} B
$$
where $\epsilon : P_\bullet \to B$ is the standard resolution
of $B$ over $A$.
\end{definition}

\noindent
Let $R$ be a ring and let $A$, $B$ be $R$-algebras. In this section we
discuss $L_{A \otimes_R B/R}$. Most of the information we want is contained
in the following diagram
\begin{equation}
\label{equation-tensor-product}
\vcenter{
\xymatrix{
L_{A/R} \otimes_A^\mathbf{L} (A \otimes_R B) \ar[r] &
L_{A \otimes_R B/B} \ar[r] &
E \\
L_{A/R} \otimes_A^\mathbf{L} (A \otimes_R B) \ar[r] \ar@{=}[u] &
L_{A \otimes_R B/R} \ar[r] \ar[u] &
L_{A \otimes_R B/A} \ar[u] \\
 &
L_{B/R} \otimes_B^\mathbf{L} (A \otimes_R B) \ar[u] \ar@{=}[r] &
L_{B/R} \otimes_B^\mathbf{L} (A \otimes_R B) \ar[u]
}
}
\end{equation}
Explanation: The middle row is the fundamental triangle
(\href{https://stacks.math.columbia.edu/tag/08QS}{equation triangle (Tag 08QS)}) for the ring maps $R \to A \to A \otimes_R B$.
The middle column is the fundamental triangle
(\href{https://stacks.math.columbia.edu/tag/08QS}{equation triangle (Tag 08QS)}) for the ring maps $R \to B \to A \otimes_R B$.
Next, $E$ is an object of $D(A \otimes_R B)$ which ``fits'' into the
upper right corner, i.e., which turns both the top row
and the right column into distinguished triangles. Such an $E$
exists by Derived Categories, Proposition \href{https://stacks.math.columbia.edu/tag/05R0}{proposition 9 (Tag 05R0)}
applied to the lower left square (with $0$ placed in the missing
spot). To be more explicit, we could for example define $E$ as the cone
(Derived Categories, Definition \ref{derived-definition-cone})
of the map of complexes
$$
L_{A/R} \otimes_A^\mathbf{L} (A \otimes_R B) \oplus
L_{B/R} \otimes_B^\mathbf{L} (A \otimes_R B)
\longrightarrow
L_{A \otimes_R B/R}
$$
and get the two maps with target $E$ by an application of TR3.
In the Tor independent case the object $E$ is zero.

\noindent
Let $A \to B$ be a ring map. In \href{https://stacks.math.columbia.edu/bibliography}{Bibliography: Lichtenbaum-Schlessinger}
there is a fairly explicit determination of $\tau_{\geq -2}L_{B/A}$
which is often used in calculations of versal deformation spaces of
singularities. The construction follows.
Choose a polynomial algebra $P$ over $A$
and a surjection $P \to B$ with kernel $I$. Choose generators
$f_t$, $t \in T$ for $I$ which induces a surjection
$F = \bigoplus_{t \in T} P \to I$ with $F$ a free $P$-module.
Let $Rel \subset F$ be the kernel of $F \to I$, in other words
$Rel$ is the set of relations among the $f_t$. Let $TrivRel \subset Rel$
be the submodule of trivial relations, i.e., the submodule of $Rel$
generated by the elements $(\ldots, f_{t'}, 0, \ldots, 0, -f_t, 0, \ldots)$.
Consider the complex of $B$-modules
\begin{equation}
\label{equation-lichtenbaum-schlessinger}
Rel/TrivRel \longrightarrow
F \otimes_P B \longrightarrow
\Omega_{P/A} \otimes_P B
\end{equation}
where the last term is placed in degree $0$. The first map is the obvious
one and the second map sends the basis element corresponding to $t \in T$
to $\text{d}f_t \otimes 1$.

\begin{definition}
\label{derived-definition-cone}
Let $\mathcal{A}$ be an additive category.
Let $f : K^\bullet \to L^\bullet$ be a morphism of
complexes of $\mathcal{A}$. The {\it cone} of $f$
is the complex $C(f)^\bullet$ given by
$C(f)^n = L^n \oplus K^{n + 1}$ and
differential
$$
d_{C(f)}^n =
\left(
\begin{matrix}
d^n_L & f^{n + 1} \\
0 & -d_K^{n + 1}
\end{matrix}
\right)
$$
It comes equipped with canonical morphisms of complexes
$i : L^\bullet \to C(f)^\bullet$ and $p : C(f)^\bullet \to K^\bullet[1]$
induced by the obvious maps $L^n \to C(f)^n \to K^{n + 1}$.
\end{definition}

\begin{lemma}
\label{lemma-tensor-product}
Let $R$ be a ring and let $A$, $B$ be $R$-algebras. The object $E$
in (\ref{equation-tensor-product}) satisfies
$$
H^i(E) =
\left\{
\begin{matrix}
0 & \text{if} & i \geq -1 \\
\text{Tor}_1^R(A, B) & \text{if} & i = -2
\end{matrix}
\right.
$$
\end{lemma}

\begin{proof}
We use the description of $E$ as the cone on
$L_{B/R} \otimes_B^\mathbf{L} (A \otimes_R B) \to L_{A \otimes_R B/A}$.
By Lemma \href{https://stacks.math.columbia.edu/tag/09CG}{lemma compare higher (Tag 09CG)} the canonical truncations
$\tau_{\geq -2}L_{B/R}$ and $\tau_{\geq -2}L_{A \otimes_R B/A}$
are computed by the Lichtenbaum-Schlessinger complex
(\ref{equation-lichtenbaum-schlessinger}).
These isomorphisms are compatible with functoriality
(Remark \href{https://stacks.math.columbia.edu/tag/09D7}{remark functoriality lichtenbaum schlessinger (Tag 09D7)}).
Thus in this proof we work with the Lichtenbaum-Schlessinger complexes.

\medskip\noindent
Choose a polynomial algebra $P$ over $R$ and a surjection $P \to B$.
Choose generators $f_t \in P$, $t \in T$ of the kernel of this surjection.
Let $Rel \subset F = \bigoplus_{t \in T} P$ be the kernel of the map
$F \to P$ which maps the basis vector corresponding to $t$ to $f_t$.
Set $P_A = A \otimes_R P$ and $F_A = A \otimes_R F = P_A \otimes_P F$.
Let $Rel_A$ be the kernel of the map $F_A \to P_A$. Using the exact sequence
$$
0 \to Rel \to F \to P \to B \to 0
$$
and standard short exact sequences for Tor we obtain an exact sequence
$$
A \otimes_R Rel \to Rel_A \to \text{Tor}_1^R(A, B) \to 0
$$
Note that $P_A \to A \otimes_R B$ is a surjection whose kernel is generated
by the elements $1 \otimes f_t$ in $P_A$. Denote $TrivRel_A \subset Rel_A$
the $P_A$-submodule generated by the elements
$(\ldots, 1 \otimes f_{t'}, 0, \ldots,
0, - 1 \otimes f_t \otimes 1, 0, \ldots)$.
Since $TrivRel \otimes_R A \to TrivRel_A$ is surjective, we find a
canonical exact sequence
$$
A \otimes_R (Rel/TrivRel) \to Rel_A/TrivRel_A \to \text{Tor}_1^R(A, B) \to 0
$$
The map of Lichtenbaum-Schlessinger complexes is given by the diagram
$$
\xymatrix{
Rel_A/TrivRel_A \ar[r] &
F_A \otimes_{P_A} (A \otimes_R B) \ar[r] &
\Omega_{P_A/A \otimes_R B} \otimes_{P_A} (A \otimes_R B) \\
Rel/TrivRel \ar[r] \ar[u]_{-2} &
F \otimes_P B \ar[r] \ar[u]_{-1} &
\Omega_{P/A} \otimes_P B \ar[u]_0
}
$$
Note that vertical maps $-1$ and $-0$ induce an isomorphism after applying
the functor $A \otimes_R - = P_A \otimes_P -$ to the source and the vertical
map $-2$ gives exactly the map whose cokernel is the desired Tor module
as we saw above.
\end{proof}
\end{document}
